\expandafter\gdef\csname SNside1\endcsname{\node[anchor=north west,inner sep=0pt,text width=48mm,font=\SideFont,align=left] at ([xshift=138mm,yshift=-466.85pt]current page.north west) {\textcolor{SNink}{\textbf{课程连接 · 目录树和进程树}}\par\vspace{1.5mm}操作系统的目录层次、进程派生关系可用树表达。硬链接若允许同一对象有多个上级入口，目录对象关系就不能直接套用``除根外唯一双亲''；先核对对象与链接的含义。};}
\expandafter\gdef\csname SNside2\endcsname{\node[anchor=north west,inner sep=0pt,text width=48mm,font=\SideFont,align=left] at ([xshift=138mm,yshift=-444.46pt]current page.north west) {\textcolor{SNink}{\textbf{最小反例 · 度为2未必是二叉树}}\par\vspace{1.5mm}最小情形只有根与一个孩子：普通无序树不区分这个孩子在左还是在右，二叉树却区分两种位置。因此``每结点至多两个孩子''尚未表达二叉树的全部结构语义。};}
\expandafter\gdef\csname SNside3\endcsname{\node[anchor=north west,inner sep=0pt,text width=48mm,font=\SideFont,align=left] at ([xshift=138mm,yshift=-95.00pt]current page.north west) {\textcolor{SNink}{\textbf{错项排除 · 叶公式只用了二叉条件}}\par\vspace{1.5mm}证明\(n_0=n_2+1\)时没有用到``完全''或``满''。若选项声称任意二叉树还必有至多一个度为1结点，就偷偷增加了完全性；一条三结点单支链即可排除。\par\smallskip\begin{tikzpicture}[x=1mm,y=1mm,>=Stealth,node distance=5mm]
\node[draw=SNink!50,fill=SNblue,rounded corners=1mm,text width=34mm,align=center,inner sep=1.2mm,font=\fontsize{7.5}{9}\selectfont] (m0) at (0,0) {A};
\node[draw=SNink!50,fill=SNblue,rounded corners=1mm,text width=34mm,align=center,inner sep=1.2mm,font=\fontsize{7.5}{9}\selectfont] (m1) at (0,-11) {B};
\node[draw=SNink!50,fill=SNblue,rounded corners=1mm,text width=34mm,align=center,inner sep=1.2mm,font=\fontsize{7.5}{9}\selectfont] (m2) at (0,-22) {C};
\draw[->,SNink] (m0) -- (m1);
\draw[->,SNink] (m1) -- (m2);
\end{tikzpicture}
\par{\fontsize{7.5}{9}\selectfont 三结点单支链有2个度1结点；叶数公式仍成立。}};}
\expandafter\gdef\csname SNside4\endcsname{\node[anchor=north west,inner sep=0pt,text width=48mm,font=\SideFont,align=left] at ([xshift=138mm,yshift=-95.00pt]current page.north west) {\textcolor{SNink}{\textbf{不变量 · 编号表达位置而非秩}}\par\vspace{1.5mm}顺序存储中下标表达满二叉树的位置，删除缺失位置后不能把剩余元素紧凑前移而继续沿用\(2i,2i+1\)。链表遍历得到的次序也不是这种位置编号。};}
\expandafter\gdef\csname SNside5\endcsname{\node[anchor=north west,inner sep=0pt,text width=48mm,font=\SideFont,align=left] at ([xshift=138mm,yshift=-95.00pt]current page.north west) {\textcolor{SNink}{\textbf{最小反例 · 先序加后序的歧义}}\par\vspace{1.5mm}最小反例为A与其唯一孩子B：B在左或在右，先序都是AB，后序都是BA。缺少中序时，序列没有给出这个单孩子的左右位置；重复值还会带来身份歧义。};}
\expandafter\gdef\csname SNside6\endcsname{\node[anchor=north west,inner sep=0pt,text width=48mm,font=\SideFont,align=left] at ([xshift=138mm,yshift=-95.00pt]current page.north west) {\textcolor{SNink}{\textbf{不变量 · 栈保存尚未完成的祖先}}\par\vspace{1.5mm}非递归中序栈中结点的左侧处理尚未完全结束，弹出后才访问根并转入右子树。后序还要记录右子树完成状态，不能把中序的弹栈时机直接挪用。};}
\expandafter\gdef\csname SNside7\endcsname{\node[anchor=north west,inner sep=0pt,text width=48mm,font=\SideFont,align=left] at ([xshift=138mm,yshift=-95.00pt]current page.north west) {\textcolor{SNink}{\textbf{错项排除 · 先看tag再沿指针}}\par\vspace{1.5mm}线索填写后，一条非空指针可能指向祖先或兄弟方向的前后继。若忽略tag把它当孩子继续递归，就可能绕回已访问位置；``指针非空''不等价于``存在子树''。};}
\expandafter\gdef\csname SNside8\endcsname{\node[anchor=north west,inner sep=0pt,text width=48mm,font=\SideFont,align=left] at ([xshift=138mm,yshift=-95.00pt]current page.north west) {\textcolor{SNink}{\textbf{不变量 · 转换不改变数据结点}}\par\vspace{1.5mm}孩子兄弟转换保留全部数据结点，但改变链接解释。原树根无兄弟，所以转换后二叉根无右孩子；森林的多个根由右链串起，不能把右链误认成原树的后代。};}
\expandafter\gdef\csname SNside9\endcsname{\node[anchor=north west,inner sep=0pt,text width=48mm,font=\SideFont,align=left] at ([xshift=138mm,yshift=-95.00pt]current page.north west) {\textcolor{SNink}{\textbf{课程连接 · 编码和外排使用同一棵树}}\par\vspace{1.5mm}哈夫曼叶权可表示字符频数，也可表示外部归并段长。每次参加合并使该叶再负担一次权值，故总合并长度等于WPL；第8章的最佳归并由此得到。};}
\expandafter\gdef\csname SNside10\endcsname{\node[anchor=north west,inner sep=0pt,text width=48mm,font=\SideFont,align=left] at ([xshift=138mm,yshift=-95.00pt]current page.north west) {\textcolor{SNink}{\textbf{错项排除 · 最优不意味着唯一}}\par\vspace{1.5mm}左右子树交换不会改变叶深和WPL；同权候选的选择也可产生不同树。应排除``不同编码必有不同平均码长''，同时避免把哈夫曼最优性扩大成完全或平衡。};}
\expandafter\gdef\csname SNside11\endcsname{\node[anchor=north west,inner sep=0pt,text width=48mm,font=\SideFont,align=left] at ([xshift=138mm,yshift=-95.00pt]current page.north west) {\textcolor{SNink}{\textbf{不变量 · 压缩路径不压缩集合}}\par\vspace{1.5mm}路径压缩把沿途结点直接挂根，保留每个元素最终代表及集合划分。根的负规模只在合并两个不同根时累加；重复合并同一集合不能再增加规模。};}
\expandafter\gdef\csname SNside12\endcsname{\node[anchor=north west,inner sep=0pt,text width=48mm,font=\SideFont,align=left] at ([xshift=138mm,yshift=-95.00pt]current page.north west) {\textcolor{SNink}{\textbf{错项排除 · 堆序没有兄弟顺序}}\par\vspace{1.5mm}小根堆的根最小，两个孩子之间却没有规定。删除根时必须先选更优孩子，再和候选比较；只因左孩子更早出现在数组，就总往左走，会破坏右侧堆序。};}
\expandafter\gdef\csname SNside13\endcsname{\node[anchor=north west,inner sep=0pt,text width=48mm,font=\SideFont,align=left] at ([xshift=138mm,yshift=-95.00pt]current page.north west) {\textcolor{SNink}{\textbf{错项排除 · 堆序没有兄弟顺序}}\par\vspace{1.5mm}小根堆的根最小，两个孩子之间却没有规定。删除根时必须先选更优孩子，再和候选比较；只因左孩子更早出现在数组，就总往左走，会破坏右侧堆序。};}
\expandafter\gdef\csname SNside14\endcsname{\node[anchor=north west,inner sep=0pt,text width=48mm,font=\SideFont,align=left] at ([xshift=138mm,yshift=-229.88pt]current page.north west) {\textcolor{SNink}{\textbf{不变量 · 压缩路径不压缩集合}}\par\vspace{1.5mm}路径压缩把沿途结点直接挂根，保留每个元素最终代表及集合划分。根的负规模只在合并两个不同根时累加；重复合并同一集合不能再增加规模。};}
\AddToShipoutPictureFG{\begin{tikzpicture}[remember picture,overlay]\ifcsname SNside\thepage\endcsname\csname SNside\thepage\endcsname\fi\end{tikzpicture}}